\documentclass[aspectratio=169]{beamer}
\setbeamersize{text margin left=0pt,text margin right=0pt}
\usetheme{Madrid}

\author{Dorota Egerlová}
\title{Kurz pedagogických dovedností}
\date{}

\usepackage[utf8]{inputenc}
\usepackage{ulem}
\usepackage{graphicx}
\usepackage{amsmath}
\usepackage{hyperref}
\usepackage{tabularx}
\usepackage{array}
\renewcommand{\arraystretch}{1.2}
\begin{document}
\begin{frame}[plain]
\title{Kurz pedagogických dovedností}
\author{Dorota Egerlová}
\maketitle



Krátký přehlední kurz zaměřený na rozvoj pedagogických kompetencí akademických pracovníků na Vysokém učení technickém v Brně. Pokrývá šest klíčových témat: vzdělávací cíle, plánování výuky, metody výuky, hodnocení a kritéria, zpětná vazba a reflexe. Kurz propojuje teoretické definice s praktickými doporučeními a aplikacemi v laboratořích a ateliérech. Cílem je podpořit přechod od teorie ke konkrétním zlepšením ve výuce.



\textbf{Dorota Egerlová}


\end{frame}
\begin{frame}[allowframebreaks]{Outline}
\tableofcontents
\end{frame}
\section{Úvod}
\begin{frame}{Úvod}
\begin{columns}[T,onlytextwidth]
\begin{column}{0.58\textwidth}
\scriptsize
\begin{itemize}

\tightlist

\item

  \textbf{Účel kurzu:} samostudijní mikrokurz pro začínající učitele -- stručné rozvíjení základních pedagogických/didaktických dovedností.

\end{itemize}


\begin{itemize}

\tightlist

\item

  \textbf{Struktura a čas:} 6 studijních slidů (doporučeně 5 min/slid) + 6 testovacích slidů (2 min/slid).

\end{itemize}


\begin{itemize}

\tightlist

\item

  \textbf{Trvání:} 6×5 + 6×2 = 42 min; maximálně 50 min včetně rezervy.

\end{itemize}


\begin{itemize}

\tightlist

\item

  \textbf{Obsah (6 témat):} vzdělávací cíle; plánování výuky; metody výuky; hodnocení a kritéria; zpětná vazba; reflexe.

\end{itemize}


\begin{itemize}

\tightlist

\item

  \textbf{Očekávané výsledky účasti:}

\end{itemize}


\begin{itemize}

\tightlist

\item

  definovat základní didaktické pojmy;

\end{itemize}


\begin{itemize}

\tightlist

\item

  rozlišit a klasifikovat hlavní prvky výuky (cíle, metody, hodnocení, zpětná vazba);

\end{itemize}


\begin{itemize}

\tightlist

\item

  analyzovat a reflektovat jejich uplatnění v konkrétní výukové situaci.

\end{itemize}


\begin{itemize}

\tightlist

\item

  \textbf{Praktické poznámky:} studijní slidy jsou vratné; testy jsou typu SBA (vyberte vždy jednu nejlepší odpověď); dodržujte doporučené časy pro optimální průběh.

\end{itemize}


\end{column}
\begin{column}{0.42\textwidth}
\centering
\includegraphics[width=\linewidth,height=0.8\textheight,keepaspectratio]{images/slide_01.png}
\end{column}
\end{columns}
\end{frame}
\section{Kontext a obsah kurzu}
\begin{frame}{Kontext a obsah kurzu}
\begin{columns}[T,onlytextwidth]
\begin{column}{0.58\textwidth}
\scriptsize
text


\begin{enumerate}

\def\labelenumi{\arabic{enumi})}

\tightlist

\item

  Definuj měřitelné cíle

\end{enumerate}


\begin{itemize}

\tightlist

\item

  Urči konkrétní výsledky a metriky úspěchu

\end{itemize}


\begin{enumerate}

\def\labelenumi{\arabic{enumi})}

\setcounter{enumi}{1}

\tightlist

\item

  Vyber obsah a metody, které k nim vedou

\end{enumerate}


\begin{itemize}

\tightlist

\item

  Zvol materiály a aktivity podporující cíle

\end{itemize}


\begin{enumerate}

\def\labelenumi{\arabic{enumi})}

\setcounter{enumi}{2}

\tightlist

\item

  Navrhni hodnocení, které ověří cíle

\end{enumerate}


\begin{itemize}

\tightlist

\item

  Použij formativní i sumativní nástroje měření

\end{itemize}


\begin{enumerate}

\def\labelenumi{\arabic{enumi})}

\setcounter{enumi}{3}

\tightlist

\item

  Přizpůsob podle kontextu a zpětné vazby

\end{enumerate}


\begin{itemize}

\tightlist

\item

  Iteruj plán na základě dat a potřeb účastníků

\end{itemize}


\begin{itemize}

\tightlist

\item

  Sleduj dlouhodobé dopady a upravuj průběžně

\end{itemize}


\end{column}
\begin{column}{0.42\textwidth}
\centering
\includegraphics[width=\linewidth,height=0.8\textheight,keepaspectratio]{images/slide_02.png}
\end{column}
\end{columns}
\end{frame}
\section{Téma 1}
\begin{frame}{Téma 1}
\begin{columns}[T,onlytextwidth]
\begin{column}{0.58\textwidth}
\scriptsize
\begin{itemize}

\tightlist

\item

  \textbf{Co jsou vzdělávací cíle:} očekávané výsledky učení; musí být studentům srozumitelně formulovány.

\end{itemize}


\begin{itemize}

\tightlist

\item

  \textbf{Srozumitelnost pro studenty:} pište je jasně, používejte běžný jazyk a sdílejte je na začátku kurzu.

\end{itemize}


\begin{itemize}

\tightlist

\item

  \textbf{Praktické doporučení:} formulujte \emph{měřitelné} cíle --- určete konkrétní výsledky a metriky úspěchu.

\end{itemize}


\begin{itemize}

\tightlist

\item

  \textbf{Použijte akční slovesa:} popište, co student ``dokáže'', ``vyjmenuje'', ``analyzuje'' apod., aby byl cíl konkrétní.

\end{itemize}


\begin{itemize}

\tightlist

\item

  \textbf{Definujte metriky:} uveďte indikátory úspěchu (např. procento, kvalita výstupu, konkrétní úroveň dovednosti).

\end{itemize}


\begin{itemize}

\tightlist

\item

  \textbf{Zarovnání s výukou:} vyberte obsah a metody, které k cílům vedou, a navrhněte hodnocení, které je ověří.

\end{itemize}


\begin{itemize}

\tightlist

\item

  \textbf{Výukové metody:} kombinujte přednášky, cvičení a praktika tak, aby podporovaly dosažení cílů.

\end{itemize}


\begin{itemize}

\tightlist

\item

  \textbf{Aplikace v praxi:} slaďte cíle s praktickými požadavky laboratoří/ateliérů a s kritérii hodnocení.

\end{itemize}


\begin{itemize}

\tightlist

\item

  \textbf{Hodnocení a rubriky:} vytvořte rubriky a zkušební úlohy přímo mapované na cíle.

\end{itemize}


\begin{itemize}

\tightlist

\item

  Doporučený čas pro tento studijní slid: \textbf{5 minut}.

\end{itemize}


\end{column}
\begin{column}{0.42\textwidth}
\centering
\includegraphics[width=\linewidth,height=0.8\textheight,keepaspectratio]{images/slide_03.png}
\end{column}
\end{columns}
\end{frame}
\section{Téma 2}
\begin{frame}{Téma 2}
\begin{columns}[T,onlytextwidth]
\begin{column}{0.58\textwidth}
\scriptsize
\begin{enumerate}

\def\labelenumi{\arabic{enumi}.}

\tightlist

\item

  Definuj dlouhodobé cíle kurzu --- zaměř se na očekávané výsledky a klíčové kompetence

\end{enumerate}


\begin{enumerate}

\def\labelenumi{\arabic{enumi}.}

\setcounter{enumi}{1}

\tightlist

\item

  Upřesni, jaké změny v dovednostech a znalostech mají účastníci dosáhnout

\end{enumerate}


\begin{enumerate}

\def\labelenumi{\arabic{enumi}.}

\setcounter{enumi}{2}

\tightlist

\item

  Rozděl je na krátkodobé cíle pro každou lekci a modul

\end{enumerate}


\begin{enumerate}

\def\labelenumi{\arabic{enumi}.}

\setcounter{enumi}{3}

\tightlist

\item

  Formuluj krátkodobé cíle tak, aby byly měřitelné a dosažitelné

\end{enumerate}


\begin{enumerate}

\def\labelenumi{\arabic{enumi}.}

\setcounter{enumi}{4}

\tightlist

\item

  Pro každou lekci: vyber obsah, metody a časovou dotaci

\end{enumerate}


\begin{enumerate}

\def\labelenumi{\arabic{enumi}.}

\setcounter{enumi}{5}

\tightlist

\item

  Plánuj konkrétní aktivity, učební materiály a časový rozpis pro lekci

\end{enumerate}


\begin{enumerate}

\def\labelenumi{\arabic{enumi}.}

\setcounter{enumi}{6}

\tightlist

\item

  Navrhni hodnocení, které ověří krátkodobé i dlouhodobé cíle

\end{enumerate}


\begin{enumerate}

\def\labelenumi{\arabic{enumi}.}

\setcounter{enumi}{7}

\tightlist

\item

  Zahrň formativní i sumativní nástroje pro kontrolu pokroku

\end{enumerate}


\begin{enumerate}

\def\labelenumi{\arabic{enumi}.}

\setcounter{enumi}{8}

\tightlist

\item

  Iteruj na základě zpětné vazby a reflexe účastníků i lektorů

\end{enumerate}


\begin{enumerate}

\def\labelenumi{\arabic{enumi}.}

\setcounter{enumi}{9}

\tightlist

\item

  Aktualizuj plán, metody a materiály podle získaných poznatků

\end{enumerate}


\end{column}
\begin{column}{0.42\textwidth}
\centering
\includegraphics[width=\linewidth,height=0.8\textheight,keepaspectratio]{images/slide_04.png}
\end{column}
\end{columns}
\end{frame}
\section{Téma 3}
\begin{frame}{Téma 3}
\begin{columns}[T,onlytextwidth]
\begin{column}{0.58\textwidth}
\scriptsize
\begin{itemize}

\tightlist

\item

  Metody výuky: způsoby zprostředkování učiva studentům (výklad, diskuse, aktivizující přístupy).

\end{itemize}


\begin{itemize}

\tightlist

\item

  Postup výběru metody --- jednoduchý algoritmus:

\end{itemize}


\begin{enumerate}

\def\labelenumi{\arabic{enumi}.}

\tightlist

\item

  Definujte výukové cíle pro danou lekci.

\end{enumerate}


\begin{enumerate}

\def\labelenumi{\arabic{enumi}.}

\setcounter{enumi}{1}

\tightlist

\item

  Zhodnoťte obsah a charakteristiky studentů.

\end{enumerate}


\begin{enumerate}

\def\labelenumi{\arabic{enumi}.}

\setcounter{enumi}{2}

\tightlist

\item

  Zvolte a kombinujte metody tak, aby cíle, obsah a studenti byly sladěni.

\end{enumerate}


\begin{enumerate}

\def\labelenumi{\arabic{enumi}.}

\setcounter{enumi}{3}

\tightlist

\item

  Ověřte účinnost a upravte (iterace podle zpětné vazby).

\end{enumerate}


\begin{itemize}

\tightlist

\item

  Proč kombinovat metody: přizpůsobení různým typům obsahu, vyšší pozornost a lepší zapojení studentů.

\end{itemize}


\begin{itemize}

\tightlist

\item

  Aktivizující přístupy: umožňují studentům aktivně pracovat s učivem --- vhodné zejména pro praktické části.

\end{itemize}


\begin{itemize}

\tightlist

\item

  Praktická doporučení pro laboratoře/ateliéry: slučujte krátké instruktážní bloky (výklad) s pracovním cvičením a stručnými diskusemi o postupu; měřte pokrok formativně.

\end{itemize}


\end{column}
\begin{column}{0.42\textwidth}
\centering
\includegraphics[width=\linewidth,height=0.8\textheight,keepaspectratio]{images/slide_05.png}
\end{column}
\end{columns}
\end{frame}
\section{Téma 4}
\begin{frame}{Téma 4}
\begin{columns}[T,onlytextwidth]
\begin{column}{0.58\textwidth}
\scriptsize
\begin{enumerate}

\def\labelenumi{\arabic{enumi})}

\tightlist

\item

  Určete, co budete hodnotit.

\end{enumerate}


\begin{itemize}

\tightlist

\item

  Zaměřte se na konkrétní dovednosti, znalosti a produkty.

\end{itemize}


\begin{enumerate}

\def\labelenumi{\arabic{enumi})}

\setcounter{enumi}{1}

\tightlist

\item

  Sdělte kritéria studentům před zadáním úkolu.

\end{enumerate}


\begin{itemize}

\tightlist

\item

  Ujistěte se, že rozumějí očekáváním a stupnicím.

\end{itemize}


\begin{enumerate}

\def\labelenumi{\arabic{enumi})}

\setcounter{enumi}{2}

\tightlist

\item

  Navážete každé kritérium na konkrétní výukový cíl.

\end{enumerate}


\begin{itemize}

\tightlist

\item

  Spojte hodnocení s tím, co má student umět po výuce.

\end{itemize}


\begin{enumerate}

\def\labelenumi{\arabic{enumi})}

\setcounter{enumi}{3}

\tightlist

\item

  Použijte jednoduchou hodnoticí rubriku nebo jasná pravidla.

\end{enumerate}


\begin{itemize}

\tightlist

\item

  Preferujte přehlednost a konzistenci při hodnocení.

\end{itemize}


\begin{itemize}

\tightlist

\item

  Poskytněte příklady úspěšných prací.

\end{itemize}


\end{column}
\begin{column}{0.42\textwidth}
\centering
\includegraphics[width=\linewidth,height=0.8\textheight,keepaspectratio]{images/slide_06.png}
\end{column}
\end{columns}
\end{frame}
\section{Téma 5}
\begin{frame}{Téma 5}
\begin{columns}[T,onlytextwidth]
\begin{column}{0.58\textwidth}
\scriptsize
Zpětná vazba jako nástroj zlepšení: stručně a akčně


\begin{itemize}

\tightlist

\item

  Hlavní rozdíl: \textbf{bodové/sumativní hodnocení} shrnuje výkon na konci; \textbf{formativní} zpětná vazba podporuje průběžné učení a zlepšování výkonu.

\end{itemize}


\begin{itemize}

\tightlist

\item

  Formativní přístup je průběžný, zaměřený na změnu chování a růst studenta.

\end{itemize}


\begin{itemize}

\tightlist

\item

  Kvalitní zpětná vazba --- buďte \textbf{konkrétní}, \textbf{včasní}, \textbf{konstruktivní}; zaměřte se na to, jak student může svou práci vylepšit.

\end{itemize}


\begin{itemize}

\tightlist

\item

  Uveďte pozorovatelné chování, konkrétní návrhy a očekávaný výsledek.

\end{itemize}


\begin{itemize}

\tightlist

\item

  Jednoduchý 3‑krokový postup pro praxi: 1) popište pozorovatelný výkon; 2) uveďte, co bylo dobré + konkrétní návrh zlepšení; 3) dejte okamžitou reakci a požádejte o následnou úpravu.

\end{itemize}


\begin{itemize}

\tightlist

\item

  Vyžadujte reflexi nebo opravu, abyste podpořili učení a sledovali pokrok.

\end{itemize}


\begin{itemize}

\tightlist

\item

  Doporučení pro laboratoře/ateliéry: krátké komentáře ihned po výkonu; ukažte příklady úspěšných řešení; pravidelně měřte pokrok formativně a dejte studentům prostor k reakci.

\end{itemize}


\begin{itemize}

\tightlist

\item

  Krátké ukázky a rychlá zpětná vazba vedou k rychlejším zlepšením.

\end{itemize}


\end{column}
\begin{column}{0.42\textwidth}
\centering
\includegraphics[width=\linewidth,height=0.8\textheight,keepaspectratio]{images/slide_07.png}
\end{column}
\end{columns}
\end{frame}
\section{Téma 6}
\begin{frame}{Téma 6}
\begin{columns}[T,onlytextwidth]
\begin{column}{0.58\textwidth}
\scriptsize
pseudo


REFLECT:


\begin{enumerate}

\def\labelenumi{\arabic{enumi}.}

\tightlist

\item

  Zaznamenej: průběh výuky, klíčové momenty, chování studentů, dosažené výsledky

\end{enumerate}


\begin{enumerate}

\def\labelenumi{\arabic{enumi}.}

\setcounter{enumi}{1}

\tightlist

\item

  Dokumentuj: konkrétní příklady problémů, úspěchů a netypických situací

\end{enumerate}


\begin{enumerate}

\def\labelenumi{\arabic{enumi}.}

\setcounter{enumi}{2}

\tightlist

\item

  Analyzuj: co fungovalo a proč; co nefungovalo a proč

\end{enumerate}


\begin{enumerate}

\def\labelenumi{\arabic{enumi}.}

\setcounter{enumi}{3}

\tightlist

\item

  Hledej vzorce v chování a výsledcích napříč hodinami

\end{enumerate}


\begin{enumerate}

\def\labelenumi{\arabic{enumi}.}

\setcounter{enumi}{4}

\tightlist

\item

  Využij: zpětnou vazbu studentů jako podklad pro rozhodování

\end{enumerate}


\begin{enumerate}

\def\labelenumi{\arabic{enumi}.}

\setcounter{enumi}{5}

\tightlist

\item

  Konzultuj s kolegy: sdílej poznatky a získej další pohledy

\end{enumerate}


\begin{enumerate}

\def\labelenumi{\arabic{enumi}.}

\setcounter{enumi}{6}

\tightlist

\item

  Upravit: metody, tempo, zadání podle analýzy

\end{enumerate}


\begin{enumerate}

\def\labelenumi{\arabic{enumi}.}

\setcounter{enumi}{7}

\tightlist

\item

  Naplánuj: konkrétní kroky, termíny a kritéria úspěchu pro změny

\end{enumerate}


\end{column}
\begin{column}{0.42\textwidth}
\centering
\includegraphics[width=\linewidth,height=0.8\textheight,keepaspectratio]{images/slide_08.png}
\end{column}
\end{columns}
\end{frame}
\section{Aplikace do praxe}
\begin{frame}{Aplikace do praxe}
\begin{columns}[T,onlytextwidth]
\begin{column}{0.58\textwidth}
\scriptsize
for každá\_lekce:


plánuj(cíle, úkoly, kritéria, časový\_rámec)


připrav\_materiály(diferenciace, pomůcky)


realizuj(metody, zapojení\_studentů)


facilituj\_diskusi\_a\_aktivity


měř(progres(formativně), výsledky(sumativně))


sbírej\_data(portfolio, testy, pozorování)


dej\_zpětnou\_vazbu(rychle, konkrétně, motivačně)


reflektuj\_a\_uprav(plán, metody, hodnocení)


dokumentuj\_a\_sdílej\_změny


\end{column}
\begin{column}{0.42\textwidth}
\centering
\includegraphics[width=\linewidth,height=0.8\textheight,keepaspectratio]{images/slide_09.png}
\end{column}
\end{columns}
\end{frame}
\section{Závěr}
\begin{frame}{Závěr}
\begin{columns}[T,onlytextwidth]
\begin{column}{0.58\textwidth}
\scriptsize
\begin{itemize}

\tightlist

\item

  Stručné shrnutí: mikrokurz obsahuje 6 klíčových témat --- vzdělávací cíle, plánování, metody, hodnocení, zpětná vazba, reflexe.

\end{itemize}


\begin{itemize}

\tightlist

\item

  Cíl mikrokurzu: definovat základní didaktické pojmy, rozlišit hlavní prvky výuky a analyzovat jejich uplatnění v konkrétních výukových situacích.

\end{itemize}


\begin{itemize}

\tightlist

\item

  Časový rámec (doporučeno): studijní část 6 × 5 min = 30 min; testovací část 6 × 2 min ≈ 12 min.

\end{itemize}


\begin{itemize}

\tightlist

\item

  Celkem doporučeně do 50 min.

\end{itemize}


\begin{itemize}

\tightlist

\item

  Praktický 4‑krokový postup (pseudokód):

\end{itemize}


planuj(cíle, úkoly, kritéria, čas) → realizuj(metody, zapojení\_studentů) → měř(formativně/sumativně) → reflektuj\_a\_uprav(plán, metody, hodnocení)


\begin{itemize}

\tightlist

\item

  Krok 1 --- plánuj: stanovte jasné cíle, konkrétní úkoly a hodnoticí kritéria spolu s časovým rozpisem.

\end{itemize}


\begin{itemize}

\tightlist

\item

  Krok 2 --- realizuj: vyberte vhodné metody a aktivizujte studenty, zajistěte zapojení a průběžnou podporu.

\end{itemize}


\begin{itemize}

\tightlist

\item

  Krok 3 --- měř: sbírejte data formativně i sumativně, hodnotte naplnění cílů a průběh činností.

\end{itemize}


\begin{itemize}

\tightlist

\item

  Reflexe jako klíč k profesnímu růstu: systematicky analyzujte průběh, posuďte naplnění cílů a plánujte změny.

\end{itemize}


\end{column}
\begin{column}{0.42\textwidth}
\centering
\includegraphics[width=\linewidth,height=0.8\textheight,keepaspectratio]{images/slide_10.png}
\end{column}
\end{columns}
\end{frame}
\end{document}
